\documentclass[11pt,a4paper]{article}
\usepackage[utf8]{inputenc}
\usepackage[T1]{fontenc}
\usepackage[brazil]{babel}
\usepackage[margin=2.5cm]{geometry}
\usepackage{booktabs}
\usepackage{amsmath,amssymb}
\usepackage{caption}
\usepackage{setspace}
\usepackage{enumitem}
\usepackage{placeins}
\usepackage{longtable}
\usepackage{float}
\usepackage{xcolor}
\usepackage[most]{tcolorbox}
\usepackage[hidelinks]{hyperref}
\setstretch{1.35}
\captionsetup{font=small,labelsep=colon}
\setcounter{tocdepth}{2}

\newtcolorbox{licao}[1]{colback=gray!8,colframe=gray!60,fonttitle=\bfseries,title={#1},breakable}

\title{Exploração de sinal de lucro a partir de holdings de fundos brasileiros\\
\large Documento de trabalho --- NÃO faz parte do TCC}
\author{João Paulo Zangrandi\\
Fundação Getulio Vargas --- Escola de Economia de São Paulo (FGV--EESP)}
\date{15 de agosto de 2026}

\begin{document}
\maketitle

\begin{center}
\fbox{\parbox{0.9\textwidth}{\centering\small
\textbf{Este documento é exploratório, não é parte do TCC oficial
(\texttt{TCC\_finalV2.tex}).} Registra, com total transparência, o processo
completo de busca por um sinal de lucro a partir dos dados que o TCC já
constrói --- mais de 700 especificações, quase 90 candidatos --- inclusive
os rejeitados, e inclusive os que quase pareceram funcionar e não
funcionavam de verdade. Nenhum conteúdo daqui deve entrar no TCC sem antes
ser discutido com o orientador.}}
\end{center}

\tableofcontents
\newpage

% ======================================================================
\section{Motivação e como ler este documento}

O TCC mostra que dá para prever razoavelmente bem o \emph{peso} que um
fundo brasileiro vai ter numa ação no mês seguinte. A pergunta natural
depois disso é: será que essa previsibilidade serve para \textbf{ganhar
dinheiro}? Prever o comportamento de um fundo não é a mesma coisa que
prever o preço da ação que ele compra --- são perguntas relacionadas, mas
diferentes, e a Seção 6 do TCC já mostrava um primeiro sinal de que a
resposta pode ser não (o fluxo previsto dos fundos não prediz retorno da
ação com $R^2$ acima de 0,2\%).

Este documento é o registro completo de uma investigação muito mais ampla
dessa pergunta: peguei tudo que os dados de holdings + preços permitem
construir, pesquisei a literatura acadêmica de ponta a ponta atrás de
mecanismos conhecidos que ligam posse institucional a retorno ou risco
futuro, e testei cada um com o mesmo rigor usado no resto do TCC (fora da
amostra, corte fixo treino/teste). O critério de honestidade é simples:
\textbf{nada aqui vira ``achado'' sem sobreviver ao teste fora da amostra},
e \textbf{todo} teste é registrado, com ou sem sucesso --- inclusive os
casos em que um resultado parecia bom de cara e, investigando mais fundo,
se revelou um artefato estatístico.

\subsection{Por que registrar até o que não deu certo}

Se eu testar 100 ideias diferentes e reportar só a única que deu
$p<0{,}05$, isso não prova nada --- é exatamente o que se espera por puro
acaso (a 5\% de significância, 1 em cada 20 testes ``passa'' sem que exista
efeito nenhum). É a armadilha estatística chamada \emph{data
dredging}/\emph{p-hacking}. A defesa contra isso é registrar \textbf{tudo}
e aplicar uma régua de significância mais dura conforme o número de testes
cresce (correção de Bonferroni: com $N$ testes já feitos, o limiar vira
$0{,}05/N$ em vez de $0{,}05$). Ao longo deste documento, o limiar de
Bonferroni sobe de $0{,}05/250$ para $0{,}05/700$ conforme a contagem
avança --- registrado em cada seção.

% ======================================================================
\section{Metodologia comum a todo teste}
\label{sec:metodologia}

Pense assim: se eu quiser saber se ``ações com característica X tendem a
subir mais no mês seguinte'', a forma ingênua de testar seria juntar todas
as observações ativo-mês numa pilha só e rodar um teste-t simples. O
problema é que ações do \textbf{mesmo mês} sobem e descem juntas (todo o
mercado reage aos mesmos eventos macro naquele mês) --- então essa pilha
não tem, de verdade, tantas observações independentes quanto parece.
Tratar 500 ações em 1 mês como 500 observações independentes é como jogar
uma moeda 500 vezes mas anotar só o resultado do primeiro lançamento e
repetir 500 vezes: parece muita informação, mas na prática você só tem 1
lançamento real.

A correção correta (\textbf{Fama-MacBeth}) é: a cada mês, calcular 1 único
número resumindo o efeito daquele mês (o spread de retorno entre o
quintil mais alto e o mais baixo da característica, por exemplo). No fim,
sobra uma série temporal --- 1 número por mês, não 1 número por ação-mês
--- e o teste de significância é feito \textbf{nessa série}, com o
erro-padrão vindo da própria variação mês a mês. Isso é usado em
\textbf{todo} teste deste documento, sem exceção, depois de um episódio
(Seção~\ref{sec:comomentum}) mostrar concretamente o que acontece quando
não se faz assim.

Junto disso, três regras foram tratadas como inegociáveis do início ao fim:

\begin{licao}{As 3 regras inegociáveis (violá-las já produziu 2 quase-falsos-positivos nesta exploração)}
\begin{enumerate}
\item Significância sempre via Fama-MacBeth de verdade (acima). Nunca teste-t pooled tratando ativo-mês como observação independente.
\item Quintis/decis para carteiras long-short sempre recortados \textbf{a cada mês} (não com um corte fixo calculado no treino e aplicado ao teste) --- e sempre reportar o número mínimo de ativos por grupo-mês antes de confiar em qualquer resultado.
\item Disciplina fora da amostra idêntica ao resto do TCC: treino $<$jan/2020, teste $\geq$jan/2020.
\end{enumerate}
\end{licao}

% ======================================================================
\section{Parte I --- A busca por um sinal de RETORNO}
\label{sec:retorno}

\subsection{Visão geral: 15 candidatos clássicos, 15 rejeições limpas}

A primeira rodada testou a família de mecanismos mais estabelecida na
literatura de finanças institucionais --- fluxo mecânico de fundos, herding,
demanda revelada, breadth of ownership, demanda latente (resíduo do
próprio modelo de peso da Etapa~1 do TCC) --- pesquisando a literatura
clássica (Coval-Stafford 2007, Lou 2012, Lakonishok-Shleifer-Vishny 1992,
Chen-Hong-Stein 2002, Koijen-Yogo 2019) e a mais recente disponível
(McLean-Pontiff-Reilly 2025, Li-Sokolinski-Tamoni 2026, Chincarini et al.\
2026). Nenhum sobreviveu ao teste fora da amostra.

\begin{table}[H]\centering\small
\caption{Resumo dos 12 candidatos clássicos rejeitados (retorno)}
\begin{tabular}{@{}p{5.2cm}p{7.5cm}@{}}
\toprule
Candidato & Resultado \\
\midrule
Flow-Induced Trading (Lou 2012) & $R^2_{\text{OOS}}$ negativo em 3 de 5 horizontes; $\beta$ instável, chega a trocar de sinal \\
Herding (LSV 1992) & $R^2_{\text{OOS}}$ entre $-0{,}015\%$ e $0{,}088\%$ \\
Demanda agregada revelada & $R^2_{\text{OOS}}$ entre $-0{,}08\%$ e $0{,}18\%$ \\
Fire sale extremo (Coval-Stafford) & $R^2_{\text{OOS}}$ negativo em 3 de 4 horizontes \\
Sinal dos fundos-líder & $R^2_{\text{OOS}}$ entre $-0{,}001\%$ e $0{,}011\%$ \\
Breadth of Ownership (CHS 2002) & $R^2_{\text{OOS}}$ negativo em todos; $p_{\text{FM}}>0{,}46$ \\
Latent Demand (Koijen-Yogo) & melhor caso $p_{\text{FM}}=0{,}078$ \\
Líderes ponderados por skill & melhor caso $p_{\text{FM}}=0{,}059$ \\
Restrição a small caps (3 sinais) & todos $p_{\text{FM}}$ entre $0{,}22$ e $0{,}82$ \\
E[FIT] (fluxo esperado, Lou fiel) & $R^2_{\text{OOS}}$ positivo mas $p_{\text{FM}}$ entre $0{,}21$ e $0{,}95$ \\
FIT suavizado (média móvel) & $h=3$ isolado passa ($p=0{,}048$), inconsistente com vizinhos \\
Ownership $\times$ Momentum & sinal oposto ao hipotetizado, inconsistente entre horizontes \\
\bottomrule
\end{tabular}
\end{table}

\subsection{Aprofundamento: a saga do CIO Peer Momentum}
\label{sec:cio}

De \textbf{um} candidato entre dezenas, um se destacou --- e sua história
é o fio condutor mais importante de todo este documento, porque passou por
três fases muito diferentes: pareceu promissor, quase virou um
falso-positivo por um erro metodológico sutil, foi corrigido e ficou
\emph{mais} forte, e morreu no fim por um motivo totalmente diferente
(iliquidez), não por falta de significância estatística.

\subsubsection{O mecanismo}

Ying (2024, \emph{Journal of Financial Economics}) documenta que ações com
donos institucionais em comum exibem difusão \textbf{gradual} de
informação: quando um fundo reage a uma notícia sobre a ação A, ele
demora a reajustar posições nas \textbf{outras} ações (B, C, D\ldots) que
também carrega --- então o retorno da ``vizinhança de ownership'' de uma
ação hoje deveria prever o retorno da própria ação no mês seguinte. A
vizinhança é medida por similaridade de cosseno entre os vetores de posição
em R\$ de cada fundo (quanto mais parecida a carteira de fundos que possuem
duas ações, mais ``vizinhas'' elas são).

\subsubsection{Fase 1: o achado inicial, promissor com ressalvas}

O primeiro teste (candidato \#15, metodologia inicial) achou, em $h=1$
(retorno do mês seguinte): $R^2_{\text{OOS}}=+1{,}55\%$ (o maior de toda a
busca até então), Fama-MacBeth $t=2{,}51$, $p=0{,}021$. Passou em quase
todas as checagens de robustez da época --- amostra saudável, sobrevive a
controlar pelo retorno próprio da ação, sobrevive a 3 controles
simultâneos (tamanho, beta, retorno próprio: $t=2{,}34$, $p=0{,}028$),
sobrevive a custo de transação agressivo. Mas tinha duas rachaduras: a
magnitude bruta ($4{,}15$pp/mês, $\sim$63\%/ano) era grande demais para
ser plausível como estratégia real e persistente, e excluindo o ano de
2020 (deixando só 2021), o resultado enfraquecia para $p=0{,}057$ --- não
passava mais de 5\%. Foi registrado como ``candidato promissor, não achado
pronto'' e o documento anterior a este parou por aqui.

\subsubsection{Fase 2: quase um segundo falso-positivo, pego a tempo}

Ao retomar a investigação para tentar transformar esse sinal em algo
monetizável, reconstruí a carteira long-short do CIO Peer Momentum como
uma série de retorno de portfólio de verdade (não só um coeficiente de
regressão) --- e o resultado bateu diferente do esperado: Sharpe muito
mais alto do que fazia sentido crer de cara. Isso acendeu um alerta:
por que tão diferente do teste original?

A resposta, ao comparar as duas metodologias lado a lado na mesma base de
dados, foi desconfortavelmente familiar: o teste original (candidato
\#15/16) formava os quintis de long-short usando \textbf{breaks fixos
calculados no período de treino}, aplicados ao período de teste. Isso
parece razoável à primeira vista, mas quebra quando a distribuição do
sinal muda de regime entre treino e teste (o que aconteceu, dado que 2020
foi um ano de volatilidade atípica): em vários meses do teste, quase
\textbf{todas} as ações caíam no mesmo quintil, sobrando só \textbf{1
única ação} na perna oposta.

\begin{table}[H]\centering\small
\caption{Método antigo (breaks fixos do treino): exemplos de quintis degenerados}
\begin{tabular}{@{}lrr@{}}
\toprule
Mês & N na perna ``curta'' & N na perna ``longa'' \\
\midrule
fev/2020 & 257 & \textbf{1} \\
abr/2020 & \textbf{1} & 242 \\
out/2021 & 336 & \textbf{1} \\
\bottomrule
\end{tabular}
\end{table}

Isso é \textbf{exatamente} o mesmo tipo de artefato que já havia produzido
um falso-positivo com o candidato Comomentum
(Seção~\ref{sec:comomentum}) --- só que ninguém tinha percebido, porque o
teste original só olhava o $t$-estatístico agregado, nunca o $N$ por
grupo-mês.

\subsubsection{A correção: quintil recortado a cada mês}

Repetindo o teste com quintis recalculados a cada mês (padrão
Fama-French, a prática correta), o resultado mudou de forma reveladora: os
grupos ficaram sempre saudáveis (\textbf{mínimo de 53 ações} por perna,
nunca degenerado), e o sinal \textbf{ficou mais forte}, não mais fraco.

\begin{table}[H]\centering\small
\caption{CIO Peer Momentum: método antigo vs.\ corrigido, $h=1$}
\begin{tabular}{@{}lrrrr@{}}
\toprule
Métrica & Método antigo (breaks fixos) & Método corrigido (recorte mensal) \\
\midrule
$N$ mínimo por perna & 1 (degenerado) & 53 \\
Sharpe (24 meses) & 1,895 (inflado por artefato) & 1,531 \\
$p$ (24 meses) & 0,021 & 0,041 \\
$p$ (só 2021, sem COVID) & \textbf{0,057} (não passa 5\%) & \textbf{0,026} (mais forte!) \\
\bottomrule
\end{tabular}
\end{table}

O padrão que antes parecia uma fraqueza --- o resultado sumir sem 2020 ---
era, na verdade, um artefato da metodologia de corte fixo interagindo mal
justamente com o ano atípico. Corrigido, o sinal sobrevive: Fama-MacBeth
com controles de tamanho, beta e retorno próprio dá $t=2{,}22$,
$p=0{,}037$; $h=3$ ($t=1{,}48$, $p=0{,}154$) e $h=6$ ($t=1{,}76$,
$p=0{,}095$) continuam mais fracos, consistente com a teoria de que a
difusão de informação se resolve rápido.

\begin{licao}{Lição \#2: um resultado que ``enfraquece sem 2020'' pode ser o sintoma de um bug, não de fragilidade genuína}
Antes de aceitar que um sinal só existe ``por causa da COVID'', verifique
se o teste continua metodologicamente correto (grupos saudáveis) dentro e
fora dessa janela. Aqui, o problema não era o período --- era que o corte
fixo de quintil, calibrado no treino, ficava mal calibrado justamente no
período de maior mudança de regime.
\end{licao}

\subsubsection{Fase 3: a morte por iliquidez, não por falta de significância}

Com o sinal estatisticamente mais sólido do que nunca, o passo óbvio era
testar se dava para transformá-lo em estratégia. Duas checagens mataram
essa esperança:

\textbf{Custo de transação.} Com apenas 100 pontos-base de custo
\emph{round-trip} por perna --- não é um custo agressivo --- o resultado
líquido já fica \textbf{negativo} ($-4{,}2\%$/ano). Só sobra algo
realmente positivo com custo muito baixo (30bps, $+12{,}6\%$/ano),
irrealista para o tipo de papel envolvido.

\textbf{Composição da perna vendida.} Os tickers mais frequentes na perna
que a estratégia manda vender são: CEEB3, CEED3, DOHL4, CEED4, KLBN4,
MDNE3, CSRN3, HBOR3, RAPT3, TRIS3, CARD3, CMIG4, EVEN3, GFSA3 --- em boa
parte, small/micro caps com aluguel de ação (necessário para venda a
descoberto) provavelmente caro demais ou simplesmente indisponível no
Brasil.

Restringindo o universo a ações realmente líquidas (top 1/3 a 2/3 por
posse institucional total) \textbf{desde o início}, os dois testes que
restam contam a história completa:

\begin{table}[H]\centering\small
\caption{CIO Peer Momentum, universo restrito a ações líquidas}
\begin{tabular}{@{}p{6.5cm}rr@{}}
\toprule
Versão & Resultado & Interpretação \\
\midrule
Long-only (só comprar o quintil top, universo líquido) &
excesso $\approx 0$, $t\approx 0{,}02$--$0{,}21$, $p>0{,}8$ &
sem valor nenhum \\
Long-short (universo líquido top 33\%) &
Sharpe $=1{,}04$, $t=1{,}48$, $p=0{,}154$ &
perde significância \\
Long-short (universo líquido top 50\%) &
Sharpe $=1{,}26$, $t=1{,}79$, $p=0{,}087$ &
perde significância \\
\bottomrule
\end{tabular}
\end{table}

A conclusão mais importante desses números: \textbf{o lucro inteiro do
long-short vinha da perna vendida em papéis ilíquidos}. A perna comprada
sozinha --- a única parte que um fundo tradicional (majoritariamente
\emph{long-only}, como os que o TCC estuda) conseguiria de fato
implementar --- não tem valor algum.

\begin{licao}{Lição \#3: ``achado estatisticamente real'' $\neq$ ``estratégia implementável''}
O CIO Peer Momentum é, tecnicamente, o resultado mais limpo e robusto de
toda a busca por sinal de retorno. E mesmo assim não é um sinal de lucro
tradeable --- porque a fonte do retorno é uma posição que a maioria dos
fundos reais não consegue montar. Duas perguntas diferentes: ``isso é
real?'' e ``dá para lucrar com isso?''.
\end{licao}

\subsection{Resumo dos $\sim$75 candidatos adicionais de retorno testados}

Depois da correção do CIO Peer Momentum, uma rodada final testou
praticamente todo ângulo metodológico razoável ainda não coberto,
explicitamente relaxando a exigência de tradeabilidade (o objetivo virou
achar \emph{qualquer} sinal estatisticamente defensável, mesmo sem ser
implementável na prática). Seis agentes de pesquisa trabalharam em
paralelo, cada um cobrindo uma família de mecanismos, todos seguindo as 3
regras inegociáveis da Seção~\ref{sec:metodologia}.

\begin{longtable}{@{}p{5.5cm}p{7.2cm}@{}}
\caption{Resumo compacto de todos os ângulos adicionais testados (retorno)}\\
\toprule
Família de sinal & Resultado \\
\midrule
\endhead
Anomalias clássicas (52-week high, reversão curto prazo, idio-vol) & nulo isoladamente; melhor interação com ownership $p=0{,}0048$, não bate Bonferroni \\
52-week-high $\times$ HHI (double-sort) & inverteu de direção vs.\ interação com tamanho de posse; nenhuma versão bate Bonferroni \\
Comomentum (ver Seção~\ref{sec:comomentum}) & falso-positivo identificado e descartado \\
Stock IQ (Cao-Fang-Lee) & nulo em todas as especificações \\
Ensemble de Machine Learning (candidato original) & nulo/fraco \\
Grid search sistemático (108 especificações, CIO) & confirma $h=1$ robusto, $h=3$ definitivamente nulo \\
Window dressing / calendário & nulo \\
Persistência de performance de FUNDO & rejeitada (44 especificações) \\
Low volatility (Brasil) & nulo isolado e interagido com ownership/CIO \\
CIO $\times$ volatilidade prevista & interação nula ($t=0{,}07$) \\
Combinação de 3 sinais (CIO+52wk-high+reversão) & melhor resultado da rodada, $p=0{,}030$, ainda longe de Bonferroni \\
Choque/atenção via variação de posição & nulo limpo \\
CIO $\times$ beta do fundo (segmentação por convicção) & inconsistente entre horizontes, ruído \\
Best ideas / convicção crescente (Cohen-Polk-Silli) & nulo, poder estatístico baixo \\
Ranking composto multi-sinal (3 métodos) & todos nulos; combinar sinais fracos piora o único sinal real \\
Centralidade de rede de posse (grau, autovetor) & achado promissor revelado como confound de tamanho ($r=0{,}67$--$0{,}73$) \\
Regime de vol $\times$ sinais existentes (breadth, herding, E[FIT], reversão) & 2 quase-acertos, ambos artefatos (ver Seção~\ref{sec:quase}) \\
Dispersão de crenças entre fundos (DMS 2002) & nulo/invertido; variação da dispersão não sobrevive a controles \\
Rotação setorial via fluxo institucional & nulo (proxy heurística de setor, 72\% dos tickers classificados) \\
Pairs trading via overlap de posse (CIO) & nenhuma evidência de reversão; sinal de continuação nunca robusto \\
ML walk-forward (XGBoost raso, refit semestral) & $R^2_{\text{OOS}}$ pequeno/negativo; nenhum long-short significativo \\
Outros 4 choques de mercado (retorno) & nulo, confirma o padrão da COVID (Seção~\ref{sec:vol-eventos}) \\
Tamanho médio dos fundos donos $\to$ retorno & $p$ entre 0,01 e 0,04, não bate Bonferroni, plausível efeito-tamanho disfarçado \\
\bottomrule
\end{longtable}

\textbf{Nenhum desses $\sim$75 candidatos adicionais produziu um sinal de
retorno novo, robusto e independente do CIO Peer Momentum.}

% ======================================================================
\section{Parte II --- A busca por um sinal de VOLATILIDADE}
\label{sec:vol}

Um giro importante na investigação: em vez de perguntar ``será que
holdings de fundos preveem \emph{para onde} o preço vai'', perguntar
``será que preveem \emph{o quanto} o preço vai oscilar''. É uma pergunta
diferente --- e, ao contrário de tudo na Parte I, a resposta aqui foi sim,
de forma robusta.

\subsection{O mecanismo: fragilidade de posse (Greenwood \& Thesmar 2011)}

A intuição, com uma analogia: pense numa festa lotada, mas com todo mundo
saindo por 2 portas apertadas em vez de 8 portas largas. Se algo assusta a
multidão, a saída pelas 2 portas vira um tumulto muito pior do que se
houvesse 8 saídas para distribuir a pressão. Uma ação com posse
concentrada em poucos fundos grandes é como a festa das 2 portas: se um
desses fundos precisar vender rápido (resgate, mudança de estratégia,
crise), não tem quem absorva a venda suavemente, e o preço oscila muito
mais. Uma ação espalhada entre centenas de fundos pequenos tem ``mais
portas'' --- a saída de qualquer um deles pesa pouco no total.

Greenwood \& Thesmar (2011, \emph{Journal of Financial Economics}) formalizam
essa intuição e mostram, com dados americanos, que a concentração de posse
institucional (medida por um índice tipo Herfindahl-Hirschman, HHI, sobre
a participação de cada fundo no total detido da ação) prediz volatilidade
futura --- \textbf{não} retorno futuro. É exatamente o mesmo mecanismo
teórico já usado no outro TCC do autor (\texttt{forecasting-exposure-itub4},
sobre risco de fundo-sobre-fundo), aqui replicado com um mecanismo/dado
diferente.

\subsection{Validação: da correlação simples ao benchmark rigoroso}

\begin{table}[H]\centering\small
\caption{HHI de posse $\to$ volatilidade futura, $h=12$ (Fama-MacBeth mês a mês)}
\begin{tabular}{@{}lrrr@{}}
\toprule
Especificação & $t$ & $p$ & Observação \\
\midrule
Só HHI & 8,98 & $<0{,}0001$ & 17 meses de teste válidos \\
HHI + vol.\ passada (controle) & 3,05 & 0,0076 & sobrevive ao controle mais óbvio \\
\bottomrule
\end{tabular}
\end{table}

Mas será que HHI só está reproduzindo o que a própria volatilidade passada
já diz (fato estilizado bem conhecido: vol.\ passada prediz vol.\ futura)?
O teste decisivo foi um benchmark HAR-RV (Heterogeneous Autoregressive
Realized Volatility --- usa 3 janelas de volatilidade passada, não só 1) e
comparar o ganho de $R^2$ fora da amostra ao adicionar HHI:

\begin{table}[H]\centering\small
\caption{Ganho de $R^2_{\text{OOS}}$ ao adicionar HHI sobre o benchmark HAR-RV}
\begin{tabular}{@{}lrrrr@{}}
\toprule
Horizonte & Só vol.\ 12m & HAR (3 janelas) & HAR + HHI & Ganho \\
\midrule
$h=3$  & 6,14\% & 6,86\% & 6,88\% & $+0{,}02$pp --- nada \\
$h=6$  & 8,29\% & 8,96\% & 9,36\% & $+0{,}40$pp --- pequeno \\
$h=12$ & 0,25\% & 0,20\% & 1,42\% & $+1{,}22$pp --- \textbf{$\sim$7$\times$ o benchmark} \\
\bottomrule
\end{tabular}
\end{table}

Em horizontes curtos, a volatilidade passada já explica quase tudo que HHI
explicaria. Mas em \textbf{12 meses}, a persistência mecânica de curto
prazo praticamente para de funcionar --- e é exatamente aí que a
fragilidade de posse mostra sua contribuição mais clara. Isso bate com a
leitura teórica: fragilidade é sobre risco de \textbf{liquidação
coordenada}, um evento de cauda que se manifesta mais em horizontes
longos, não uma persistência de curto prazo.

\subsection{Réplica em 4 outros choques de mercado brasileiro}
\label{sec:vol-eventos}

Um teste de alto poder estatístico repetiu o desenho ``HHI antes do
choque prevendo volatilidade depois do choque'' na crise da COVID
(fev-mar/2020, $N=217$ ações numa única regressão cross-sectional) --- e
não achou nada de conclusivo isoladamente (queda no crash: $p=0{,}108$,
com sinal na direção \textbf{oposta} à hipótese de \emph{fire sale}).
Isso motivou repetir o mesmo desenho em outros 4 episódios de choque do
mercado brasileiro, fora da COVID: Joesley Day (mai/2017), greve dos
caminhoneiros (mai/2018), e 2 choques orgânicos identificados nos próprios
dados (nov/2017, abr/2019).

\begin{table}[H]\centering\small
\caption{HHI pré-choque $\to$ volatilidade nos 6 meses seguintes, 4 eventos}
\begin{tabular}{@{}lrr@{}}
\toprule
Evento & $t$ (regressão do evento) & $p$ \\
\midrule
Joesley Day (mai/2017) & 2,76 & 0,006 \\
Greve dos caminhoneiros (mai/2018) & 1,33 & 0,180 \\
Choque orgânico (nov/2017) & 3,26 & 0,001 \\
Choque orgânico (abr/2019) & 2,49 & 0,014 \\
\midrule
\multicolumn{3}{@{}l}{Mini Fama-MacBeth entre os 4 eventos ($n=4$, $\text{df}=3$): $t=3{,}74$, $p=0{,}033$} \\
\bottomrule
\end{tabular}
\end{table}

Sinal \textbf{positivo nos 4 eventos, sem exceção} --- e o resultado fica
mais forte, não mais fraco, ao remover o maior outlier (\emph{leave-one-out}).
Isso é evidência de que o mecanismo HHI$\to$volatilidade é real e não uma
peculiaridade do episódio único da COVID, que era a limitação mais óbvia
do achado original.

\subsection{Achado novo: número de cotistas dos fundos donos}
\label{sec:cotistas}

O achado mais recente e talvez o mais surpreendente de toda a exploração
veio de olhar para um lado diferente do problema: em vez de perguntar
``quão concentrada é a posse de uma ação'' (HHI), perguntar ``quão
concentrada é a \textbf{base de clientes} dos fundos que a possuem''.

A intuição: um fundo com poucos cotistas grandes tem risco de resgate
``desigual'' --- um único cliente relevante saindo pode forçar a
liquidação de posições inteiras. Um fundo com milhares de cotistas
pequenos tem resgates que tendem a se cancelar estatisticamente (uns
entram, outros saem, o efeito líquido é suave). A variável construída ---
média (ponderada pela posição em R\$) do log(1+número de cotistas) dos
fundos donos de cada ação --- é uma medida do lado do \textbf{passivo} dos
fundos, não da concentração de posse do lado do ativo (correlação com HHI
é fraca, $-0{,}07$ a $-0{,}25$ conforme o corte).

\begin{table}[H]\centering\small
\caption{Número de cotistas dos fundos donos $\to$ volatilidade futura}
\begin{tabular}{@{}lrrr@{}}
\toprule
Horizonte & $t$ & $p$ & Observação \\
\midrule
$h=3$  & $-3{,}06$ & 0,0058 & não sobrevive a controles multivariados \\
$h=6$  & $-5{,}07$ & \textbf{0,00006} & \textbf{bate Bonferroni isoladamente} \\
$h=12$ & $-4{,}73$ & 0,00022 & muito perto de Bonferroni \\
\bottomrule
\end{tabular}
\end{table}

Direção: \textbf{mais} cotistas nos fundos donos $\to$ \textbf{menos}
volatilidade futura --- exatamente a direção que a intuição de
``resgate mais suave'' sugere. Em $h=6$ e $h=12$, o efeito sobrevive
simultaneamente a 3 controles concorrentes (HHI, tamanho do fundo dono,
tamanho da posição) \emph{e} ao controle mais difícil que existe para
volatilidade futura --- a própria volatilidade passada. Sinal consistente
em 82--86\% dos meses de teste, e continua significativo mesmo excluindo a
COVID inteira (com poucochíssima amostra remanescente, $n=6$--$10$ meses,
mas na mesma direção). É uma dimensão de informação genuinamente nova ---
não é o achado de HHI reformulado.

\begin{licao}{Ressalva honesta sobre este achado}
É um resultado de \textbf{previsibilidade de risco}, não de retorno ---
testado diretamente como preditor de retorno, deu nulo ($p$ entre 0,36 e
0,78). E, como todo o resto desta exploração, sofre da mesma limitação de
fundo: a amostra de teste (2016--2021) é curta, e mesmo um resultado
``robusto dentro do que se tem'' carrega o risco de ser sorte de período
--- o mesmo tipo de armadilha que já pegou o CIO Peer Momentum antes da
correção metodológica.
\end{licao}

% ======================================================================
\section{Parte III --- Tentando monetizar a previsão de volatilidade}
\label{sec:monetizar}

Prever volatilidade não é a mesma coisa que lucrar com essa previsão ---
volatilidade é uma medida de \emph{quanto} o preço oscila, não de
\emph{para onde}. Sem dado de opções, as formas conhecidas na literatura
de transformar previsão de risco em vantagem financeira são limitadas.
Pesquisa exaustiva (5 frentes em paralelo) mapeou o que existe:

\begin{itemize}
\item \textbf{Prêmio de variância ``de verdade'' (VRP) é definicionalmente
impossível sem opções} --- depende de uma expectativa de mercado extraída
de preço de opção (risco-neutra), que simplesmente não existe sem esse
dado. Não é uma lacuna de dado contornável.
\item \textbf{Caminho mais promissor na literatura}: usar a previsão de
volatilidade para dimensionar posição (Moreira \& Muir 2017,
\emph{Volatility-Managed Portfolios}), não para prever direção --- com um
alerta forte da literatura (Cederburg et al.\ 2020, \emph{JFE}) de que o
ganho desaparece quando reestimado genuinamente fora da amostra.
\item \textbf{Betting Against Correlation} (Asness-Frazzini-Gormsen-Pedersen 2020) --- decompõe risco em correlação vs.\ volatilidade, construído só com preço à vista.
\end{itemize}

Testados empiricamente (todos usando só o período de teste, $\geq$jan/2020):

\begin{table}[H]\centering\small
\caption{Tentativas de monetizar a previsão de volatilidade}
\begin{tabular}{@{}p{6cm}rr@{}}
\toprule
Desenho & Resultado & Interpretação \\
\midrule
Carteira \emph{tilted} por decil de HHI vs.\ equal-weight &
Sharpe $-0{,}11$ vs.\ $-0{,}04$ & pior, não melhor \\
Mercado escalado por HHI agregado vs.\ buy-hold &
Sharpe $0{,}09$ vs.\ $-0{,}04$ & nominal melhor, drawdown pior, $N=24$ \\
Vol-managed clássico (escalado por vol.\ passada, sem HHI) &
Sharpe $-0{,}91$ & catastrófico (perdeu recuperação em V pós-COVID) \\
Long-short baixo-HHI menos alto-HHI &
$p=0{,}138$, sinal invertido & não significativo \\
Betting Against Correlation (3 cortes de liquidez) &
$p$ entre 0,31 e 0,78 & nulo \\
\bottomrule
\end{tabular}
\end{table}

\textbf{Nenhum desenho produziu um sinal de lucro robusto} --- o mais
``positivo'' repousa sobre 24 observações mensais dominadas por um único
evento (COVID), poder estatístico claramente insuficiente para validar
qualquer estratégia de portfólio com confiança.

% ======================================================================
\section{As três armadilhas evitadas: uma comparação lado a lado}
\label{sec:quase}

Três vezes nesta exploração, um resultado pareceu bom demais e, ao
investigar a fundo, se revelou artefato. Vale olhar as três juntas, porque
o padrão que as une é mais instrutivo do que cada uma isolada.

\subsection{Armadilha \#1: Comomentum}
\label{sec:comomentum}

O candidato com a evidência mais forte na literatura original (Lou \&
Polk, identificação quase-causal, $\sim$9\%/ano), adaptado de retorno
semanal (o dado original) para mensal (o que temos). Primeira leitura:
$h=6$, $t=2{,}21$, $p=0{,}041$ --- parecia passar. Mas a magnitude
($+6{,}08$pp/\textbf{mês}, mais de 100\%/ano) era implausível.
Investigando: os grupos de comomentum tinham entre \textbf{1 e 13 ações
por mês} --- a ``média do grupo'' era dominada por 1--2 ativos
idiossincráticos. Um único mês contribuiu $+43{,}6$pp de spread sozinho.

\subsection{Armadilha \#2: CIO Peer Momentum (metodologia original)}

Já detalhada na Seção~\ref{sec:cio}: breaks de quintil fixos calculados no
treino, aplicados ao teste, geravam grupos de até \textbf{1 única ação} em
vários meses do período de teste --- mesmo padrão da Armadilha \#1, só que
sem ter sido percebido inicialmente porque o teste olhava só o
$t$-estatístico agregado.

\subsection{Armadilha \#3: centralidade de rede de posse}

No meio da rodada final, um sinal de centralidade (grau ponderado de cada
ação na rede de posse comum institucional) pareceu promissor ($h=6$
$p=0{,}022$, $h=12$ $p=0{,}049$). Diagnóstico revelou dois problemas
simultâneos: (i) correlação de $0{,}67$ a $0{,}73$ com o tamanho (log do
capital total detido) --- controlando por tamanho, o efeito vira
$p=0{,}20$ a $0{,}95$; (ii) 2 meses da COVID concentravam 42--54\% do
spread total, e excluindo 2020 o $t$ cai de $-2{,}50$ para $-0{,}18$.

\begin{table}[H]\centering\small
\caption{As três armadilhas, lado a lado}
\begin{tabular}{@{}p{3.3cm}p{4.5cm}p{4.5cm}@{}}
\toprule
Armadilha & Sinal do problema & Como foi pego \\
\midrule
Comomentum & magnitude implausível ($>$100\%/ano) & investigar composição de grupos pequenos \\
CIO (original) & enfraquecia sem 2020 (parecia fragilidade genuína) & comparar metodologia de corte de quintil lado a lado \\
Centralidade de rede & $p$ ``bom'' mas sem controle de tamanho & regressão com controle simultâneo + exclusão de sub-período \\
\bottomrule
\end{tabular}
\end{table}

\begin{licao}{O padrão comum às três armadilhas}
Em todas, o sintoma inicial era um $p$-valor convincente. Em todas, o
problema real só aparecia ao fazer uma das três perguntas de diagnóstico:
(1) quantas observações \emph{de verdade} sustentam cada grupo/mês? (2) a
magnitude é economicamente plausível, comparada com o que a literatura
documenta como \emph{recorde}? (3) o efeito sobrevive a excluir o
sub-período mais extremo, ou desaparece com ele? Nenhuma dessas perguntas
aparece automaticamente num teste de significância --- precisam ser feitas
à parte, sempre.
\end{licao}

% ======================================================================
\section{Conclusão geral e recomendação}

Depois de mais de 700 especificações testadas, cobrindo praticamente todo
ângulo metodológico razoável --- inclusive relaxar deliberadamente a
exigência de que o sinal fosse implementável na prática --- o quadro que
se consolidou é este:

\begin{itemize}
\item \textbf{Sinal de retorno/lucro: não encontrado.} O único candidato
com suporte estatístico genuíno (CIO Peer Momentum, $h=1$) depende quase
inteiramente de vender a descoberto ações ilíquidas --- não é, na prática,
uma estratégia implementável pelos fundos que o TCC estuda.
\item \textbf{Sinal de previsibilidade de risco: real, e cada vez mais
robusto.} Concentração de posse institucional (HHI) prediz volatilidade
futura em 12 meses, replicando Greenwood \& Thesmar (2011) e confirmado em
5 episódios de choque de mercado brasileiro diferentes (COVID + 4
outros). Um segundo preditor, independente do primeiro --- número de
cotistas dos fundos donos --- bate o limiar de significância mais
rigoroso aplicado nesta exploração inteira.
\item \textbf{Nenhuma tentativa de monetizar essa previsão de risco (sem
usar opções) produziu uma estratégia robusta}, dado o tamanho pequeno da
amostra de teste disponível (2016--2021).
\end{itemize}

\textbf{Recomendação:} o achado de \textbf{previsibilidade} (HHI e
cotistas $\to$ volatilidade futura) é um resultado real, bem fundamentado
na literatura, sem precedente conhecido de aplicação ao mercado
brasileiro --- candidato genuíno a entrar na consolidação de achados do
TCC, como contribuição de previsibilidade de risco (não de retorno). A
busca por um sinal de lucro tradeable, por outro lado, não teve sucesso
apesar do esforço exaustivo --- e esse resultado nulo, tratado com o mesmo
rigor de qualquer outro, é em si uma conclusão honesta e defensável.

\end{document}
